\subsection*{Appendix F: Realized Price Under Intermediaries}

A physical product has a consumer-facing price of ₹1,000. Suppose ₹200 is retained by a retail intermediary and ₹100 by distributors or associated trade economics, leaving ₹700 for the manufacturer before any relevant taxes or adjustments. If manufacturing and fulfillment cost ₹450 per unit, the simplified contribution before other selling and fixed costs is ₹250, \textbf{not ₹550}. Sales reports using retail price as manufacturer revenue would be misleading.

The same logic applies to education sold through partners, software sold through resellers, and marketplaces where a platform intermediates transactions. The founder must identify which economic entity is principal, which revenue belongs to whom, and how contractual concessions are accounted for. Realized-price analysis is the bridge between pricing strategy and formal unit economics.

\subsection*{Appendix G: A Negotiating Process That Does Not Rely on Pressure}

For complex B2B contracts, negotiation should begin with a shared definition of the problem and scope. Show the current alternative and value model; explain implementation obligations; disclose the base service and any optional service levels; and define a change process. If a buyer asks for a reduction, explore whether they need less scope, longer payment terms, lower risk, a different performance standard, or a smaller initial pilot. A discount with unchanged obligations should be a deliberate choice, not an automatic response to disagreement.

A founder should know the \textbf{walk-away economics}: direct cost, committed capacity, expected sales effort, credit risk, minimum fair contribution, and the opportunity cost of rejecting another customer. Yet walking away can also be justified by ethics or mission even when the transaction would be profitable. This creates a negotiating stance based on defensible terms rather than bluffing about false demand or manufactured scarcity.

\subsection*{Appendix H: Pricing Experiment Design Card}

\begin{quote}
\small
\textbf{QUESTION:} Which price or package choice will this test inform?\\
\textbf{BUYER:} How are eligible buyers identified and reached?\\
\textbf{UNIT:} What exactly is priced?\\
\textbf{ALTERNATIVE:} What might the buyer choose instead?\\
\textbf{TREATMENT:} What changes, and what stays fixed?\\
\textbf{ASSIGNMENT:} Randomized, sequential, geographic, or other?\\
\textbf{FAIRNESS:} How are comparable customers treated responsibly?\\
\textbf{PRIMARY OUTCOME:} Actual paid purchase or another hard action?\\
\textbf{SECONDARY OUTCOMES:} Refunds, activation, quality, CAC, support?\\
\textbf{POWER / UNCERTAINTY:} What can this sample plausibly distinguish?\\
\textbf{HYPOTHESIS:} What result would justify keeping the change?\\
\textbf{DOWNSIDE:} What might damage trust or economics?\\
\textbf{REVIEW:} When will the decision be made?
\end{quote}

An experiment without a decision rule can produce endless ``insights'' and no improvement. An experiment that changes everything at once can validate an offer but not isolate the cause of success. Choose the study design according to the decision you actually need to make.

\subsection*{Appendix I: A Founder's Dynamic-Pricing Decision Tree}

Start by asking whether \textbf{capacity is perishable or costs genuinely change}. If not, dynamic pricing may introduce more confusion than value.

If yes, ask whether the seller can measure available capacity and forecast demand well enough for adjusted prices to improve allocation.

Next ask whether buyers can understand and anticipate the rule.

Then model the effect on long-term customer trust, access, and regulatory exposure.

Finally test whether changing prices actually improves contribution and fulfillment, rather than merely raising revenue during scarcity.

For a small educational institution, the simplest effective mechanism may be fixed public tuition, explicit enrollment windows, and a transparent pilot fee. Sophisticated software-mediated dynamic pricing is not an achievement in itself. It is justified only if its measurable benefit exceeds the operating and trust costs it creates.

\subsection*{Appendix J: 10 Competitor-Response Questions}

\begin{enumerate}
  \item If a competitor copies our price, what remains valuable?
  \item If they undercut, are they cheaper because of lower cost or because of weaker service?
  \item If they imitate our premium promise, can we still demonstrate more credible outcomes?
  \item Are our reviews and assessments calibrated well enough to be meaningfully different?
  \item Is our sales channel protected by permissioned relationships and authentic proof?
  \item Does our package require activities that the competitor's operating model makes expensive?
  \item Can we narrow the target job further to serve it with higher precision?
  \item Can we eliminate a cost that the buyer does not value?
  \item Can we lower price without violating our core quality standard?
  \item Is a different buyer, payer, or institutional partner more appropriate?
\end{enumerate}

The point is not to construct an unassailable price. It is to become an organization that can keep earning a share of customer value as alternatives change.
